
% ==================================================
% 2. POSITIONING
% ==================================================

\section{Positioning}


% ==================================================
% 2.1 PROBLEM STATEMENT
% ==================================================

\subsection{Problem Statement}

Patients who take multiple medications, especially those
with chronic conditions, may receive prescriptions from
different doctors. Different brand names, unclear
prescriptions, and limited medication knowledge make it
difficult to identify potential risks and manage medication
schedules.

\begin{tabularx}{\textwidth}{|L{0.27\textwidth}|Y|}
\hline

\textbf{The problem of}
& Difficulty identifying medications, understanding
potential safety risks, and managing medication
schedules, especially when taking medicines
prescribed by different doctors.
\\ \hline

\textbf{Affects}
& \PAChange{Patients taking multiple medications,
especially older adults and people with chronic
conditions, and healthcare professionals who
review medication information.}
\\ \hline

\textbf{The impact of which is}
& Patients may overlook drug interactions or
duplicate active ingredients, misunderstand
incomplete safety information, or miss doses.
Healthcare professionals may receive incomplete
medication information.
\\ \hline

\textbf{A successful solution would}
& Help users enter and confirm medication information,
identify potential risks, understand safety warnings
and data limitations, seek professional review
when needed, and follow confirmed medication schedules.
\\ \hline

\end{tabularx}


% ==================================================
% 2.2 PRODUCT POSITION STATEMENT
% ==================================================

\subsection{Product Position Statement}

\begin{tabularx}{\textwidth}{|L{0.27\textwidth}|Y|}
\hline

\textbf{For}
& \PAChange{Vietnamese patients taking medications,
especially those managing multiple prescriptions.}
\\ \hline

\textbf{Who}
& Need an easy way to understand potential medication
risks, manage confirmed schedules, and know when
professional advice is needed.
\\ \hline

\textbf{The Medication Safety Assistant}
& \PAChange{Is a responsive Web research prototype
for medication safety support and adherence tracking.}
\\ \hline

\textbf{That}
& Supports prescription images, natural-language
text, voice, and manual input; checks drug interactions
and duplicate active ingredients; provides safety
warnings, simulated HCP review, and medication reminders.
\\ \hline

\textbf{Unlike}
& \PAChange{Separate medicine search tools,
drug interaction checkers, and reminder
applications used for different tasks.}
\\ \hline


\textbf{Our product}
& \PAChange{Combines Vietnamese medication
identification, safety checks, and
confirmed-schedule reminders in one
workflow. It clearly reports uncertain
results, applies safety rules to new
reminders, and does not independently
change prescribed treatment.}
\\ \hline


\end{tabularx}
